\section{Conclusion}
Examiner Coach demonstrates how a modular RAG architecture can support formative
OSCE examiner feedback training. By combining speech transcription, structured
feedback-quality modelling, evidence retrieval, and LLM-based coaching, the
system provides a practical prototype for AI-assisted medical education. The
prototype evaluation suggests that retrieval, especially HyDE-based retrieval,
can improve calibration compared with a no-RAG baseline, although further expert
validation is required.

The project's main value lies in making feedback quality computationally
inspectable without reducing it to a single opaque score. The six-criterion
rubric gives examiners concrete dimensions for reflection, while the modular RAG
pipeline provides a way to ground evaluation and coaching in educational
evidence. As a research prototype, Examiner Coach therefore offers both a
technical contribution to applied RAG design and a practical contribution to
medical education feedback training.

In practical terms, the prototype responds to a concrete institutional training
need: examiners are expected to deliver structured feedback, but require a
scalable and privacy-conscious way to practise before the examination setting.
Examiner Coach is not a replacement for expert-led examiner training, but it
offers a possible practice layer between large-group instruction and real OSCE
delivery. It also demonstrates how KISSKI-style AI service infrastructure can
support an applied health-education scenario, connecting research on AI methods
with the practical requirements of sensitive institutional environments.
